La descripción completa de cada una de las variables se encuentra en el Anexo B.
\begin{enumerate}
    \item Procesamiento
    \itemsep=0pt\topsep=0pt\partopsep=0pt
    \parskip=0pt\parsep=0pt
    \begin{itemize}
        \item Se crearon procesos o ‘pipelines’ automatizados de ingesta, transformación y carga (ETL en inglés) para obtener datos de GlobColour. 
        El proceso de ingesta utiliza el protocolo FTP al conectarse a los servidores remotos con una clave privada solicitada directamente a los administradores del proyecto por correo electrónico. El acceso es público, pero debe hacerse esa solicitud justificando su uso. 
        El proceso de transformación se encargó de convertir los archivos .nc (no legibles por humanos) en archivos .csv tabulares. Cada uno de los archivos .nc contenía información de una sola variable en una fecha en particular en varios puntos geográficos (con coordenadas). El proceso de transformación se encargó entonces de agrupar estos archivos por variable, por zona y por fecha, para que después se consolidaran en una tabla que puede exportarse como .csv. 

        \begin{figure}[htbp]
            \centering
            \includegraphics[width=1\linewidth]{Imagen1.png}
            \caption{ETL de datos de GlobColour y COBI}
            \label{fig:etl_globcolour}
        \end{figure}

        \item Para homogeneizar el ajuste y entrenamiento de los modelos, se hizo un procesamiento adicional en donde se hace un desplazamiento de 3 meses de las variables X hacia adelante. Esto con el fin de utilizar las variables X de hoy para predecir 3 meses en el futuro. El desplazamiento se decidió debido a que es un periodo de tiempo acordado con COBI para que los pescadores consideren el modelo útil.
        
        \begin{figure}[htbp]
            \centering
            \includegraphics[width=0.9\linewidth]{4DkMD.png}
            \caption{Ejemplo ilustrativo del desplazamiento para una variable, en este caso el desplazamiento es hacia atrás. En azul los datos originales y en amarillo los datos desplazados}
            \label{fig:desplazamiento}
        \end{figure}
        
        \item Para obtener los datos de los sensores de COBI, quienes los recopilan cada 6 meses, se utiliza una conexión directa a OneDrive para recuperar sus datos en formato .csv. 

        \item Una vez cargados los datos se realizó una agregación espacial y temporal, agrupando por promedio todas las observaciones de la misma fecha y de las regiones cercanas a San Quintín que corresponden a la zona de pesca de la unidad económica LITORAL DE BAJA CALIFORNIA S DE PR DE RL. Esta decisión se tomó debido a la disponibilidad de datos para esta unidad económica sobre la pesca de Langosta.
        
        Los datos finales se almacenan en una tabla donde cada observación corresponde al agregado por promedio de todas las variables explicativas para una fecha en particular en toda la zona de interés.

        \begin{figure}[htbp]
            \centering
            \includegraphics[width=0.8\linewidth]{mapa.png}
            \caption{Zona de pesca designada a la unidad económica LITORAL DE BAJA CALIFORNIA S DE PR DE RL}
            \label{fig:zona_pesca}
        \end{figure}

        Para reducir el ruido inducido por los meses de veda, donde a pesar de que las condiciones del mar pueden ser normales pero no hay pesca, decidimos reducir el conjunto de entrenamiento y prueba a sólo los días de temporada de pesca de Langosta: del 15 de septiembre de cada año al 15 de febrero del año siguiente.
        
    \end{itemize}


    \item Limitaciones
    Hay diferentes limitaciones que debemos tomar en cuenta cuando el proyecto se extienda a otras regiones y se escale:
    \begin{itemize}
        \item  A pesar de que los datos que usamos son públicos —a excepción de los de COBI, sobre los cuales tenemos derechos por convenio—, dependemos de su disponibilidad y del correcto funcionamiento de sus propios procesos de publicación. 
        \item Además, los datos oceanográficos obtenidos de COBI están limitados en cuanto a su disponibilidad temporal y geográfica. Los sensores de COBI se recuperan cada 6 meses y su uso está limitado a zonas cercanas.

        \item Una limitación más es la escasez de datos sobre arrubos pesqueros (de donde obtenemos nuestra variable objetivo): A pesar de tener más de 90 mil observaciones, al agregar los datos por especie y por unidad económica (a cada unidad económica le corresponde una zona de pesca en particular), las observaciones se reducen drásticamente. El ejercicio de este proyecto se enfoca en la especie Langosta y en la unidad económica  LITORAL DE BAJA CALIFORNIA S DE PR DE RL debido a que aquí es donde encontramos más de 400 observaciones útiles (el número máximo comparado con las demás especies y unidades económicas). En futuras implementaciones de este proyecto, esperamos contar con el apoyo de las unidades económicas para aumentar los datos disponibles.
        
        \item Por último, los arribos pesqueros en general son muy limitados y son información valiosa para cada comunidad. Esperamos que este trabajo sea un incentivo a seguir generando estos datos para mejorar los modelos de predicción en el futuro.
    \end{itemize}
\end{enumerate}
